\chapter*{Abstract}

In this thesis, simulations of \acrfull{tbl} over curved walls under \acrfull{zpg} conditions are performed to enable an isolated investigation of curvature-induced effects in the \acrshort{tbl}.  
In curved geometries, a streamwise pressure gradient naturally develops, causing pressure-gradient effects to superimpose on those induced by wall curvature complicating their independent assessment. To isolate the influence of curvature, the pressure field was manipulated to establish \acrshort{zpg} conditions throughout 
 the entire simulated domain.
This work evaluates a novel approach for constructing a curved simulation setup under \acrshort{zpg} conditions.
The method is based on the optimisation algorithm proposed by \citet{morita_applying_2022}, which performs a shape optimisation  of the upper wall of the simulated domain in order to match a prescribed distribution of the Clauser parameter $\beta$ at the lower wall.
The resulting velocity field is extracted and imposed as a \benennung{Dirichlet} boundary condition on the upper wall of a \acrlong{dns} setup.
Using this methodology,  \acrlong{dns} of \acrlong{tbl} are carried out for three convex curved geometries with different curvature characteristics.
The results are compared with simulations of curved boundary layers under reference conditions as well as with reference data from flat-plate \acrlong{tbl} simulations.
The results demonstrate that the optimisation algorithm effectively controls  the pressure field and achieves an approximately uniform\fxinfo{constant?} pressure-gradient distribution on a global scale. 
However, its performance is limited in cases with strong curvature, such as a 90° bend, where the target \acrshort{zpg} conditions cannot be maintained throughout the entire domain.

